% ==================================================================
%  WEEK 1
% ==================================================================
\chapter{Introduction, NIST Model and Cloud Architecture Basics}

\section{Overview of Computing}
The ACM Computing Curricula 2005 defines \emph{computing} as ``any goal-oriented activity
requiring, benefiting from, or creating computers''. It therefore covers designing and
building hardware and software systems, processing and managing information, scientific
studies using computers, making systems behave intelligently, creating and using
communication and entertainment media, and finding and gathering relevant information.

\subsection{Course Roadmap}
\begin{enumerate}[leftmargin=*]
  \item \textbf{Introduction:} overview of computing, NIST model, properties, characteristics,
        disadvantages, role of open standards.
  \item \textbf{Architecture:} cloud stack, service models (IaaS, PaaS, SaaS), deployment models.
  \item \textbf{Service management:} Service Level Agreements (SLAs) and cloud economics.
  \item \textbf{Resource management} in cloud computing.
  \item \textbf{Data management:} scalability, databases and data stores, large-scale data
        processing.
  \item \textbf{Cloud security:} infrastructure security, data security and storage,
        identity and access management, access control, trust, reputation, risk.
  \item \textbf{Case studies:} open-source and commercial clouds, cloud simulators.
  \item \textbf{Research trends:} fog computing and others.
\end{enumerate}

\subsection{Trends in Computing}
The lecture traces the evolution \emph{Distributed Computing $\rightarrow$ Grid Computing
$\rightarrow$ Cluster Computing $\rightarrow$ Utility Computing $\rightarrow$ Cloud Computing}.

\begin{figure}[H]
\centering
\begin{tikzpicture}[node distance=6mm]
  \node[dark,minimum width=2.5cm] (a) {Distributed};
  \node[box,minimum width=2.5cm,right=of a] (b) {Grid};
  \node[box,minimum width=2.5cm,right=of b] (c) {Cluster};
  \node[box,minimum width=2.5cm,right=of c] (d) {Utility};
  \node[dark,minimum width=2.5cm,right=of d] (e) {Cloud};
  \draw[arr] (a)--(b); \draw[arr] (b)--(c); \draw[arr] (c)--(d); \draw[arr] (d)--(e);
\end{tikzpicture}
\caption{Evolution of computing paradigms leading to cloud computing.}
\end{figure}

\section{Distributed Computing}
\begin{defbox}{Distributed System}
A distributed system consists of several \textbf{autonomous computational entities}, each
with its \textbf{own local memory}, that communicate with each other by \textbf{message passing}
(Wikipedia). In \emph{Operating System Concepts}, processors communicate through lines such as
high-speed buses or telephone lines, and each processor has its own local memory.
\emph{Distributed computing} is the field that studies distributed systems and uses them to
solve computational problems.
\end{defbox}

Early computing used a single processor (\emph{uni-processor} or \emph{centralized}
computing). Examples of distributed systems are the Internet, ATM (bank) machines,
intranets/workgroups, and the emerging landscape of ubiquitous network-connected devices.

\subsection{Computers in a Distributed System}
\begin{itemize}[leftmargin=*]
  \item \textbf{Workstations:} used by end users to perform computing.
  \item \textbf{Server systems:} provide resources and services.
  \item \textbf{Personal assistance devices:} handheld computers connected over wireless links.
\end{itemize}

\subsection{Common Properties}
\begin{itemize}[leftmargin=*]
  \item \textbf{Fault tolerance:} when one or some nodes fail, the system still works
        (only performance degrades); the status of every node must be checked.
  \item \textbf{Partial role of each node:} each computer has a limited, incomplete view of
        the system and may know only one part of the input.
  \item \textbf{Resource sharing:} users share computing power and storage.
  \item \textbf{Load sharing:} dispatching tasks to different nodes balances system load.
  \item \textbf{Easy expansion:} adding nodes should take little or no time.
  \item \textbf{Performance:} parallel computing is considered a subset of distributed computing.
\end{itemize}

\subsection{Why Distributed Computing?}
\begin{itemize}[leftmargin=*]
  \item \textbf{Nature of application.}
  \item \textbf{Performance:} \emph{computing-intensive} tasks (e.g.\ Monte Carlo
        estimation of $\pi$) and \emph{data-intensive} tasks (e.g.\ Facebook, Large Hadron
        Collider data processing).
  \item \textbf{Robustness:} no \emph{Single Point Of Failure} (SPOF); other nodes can execute
        the task of a failed node.
\end{itemize}

\subsection{Distributed Applications: Client-Server vs.\ Peer-to-Peer}
A distributed application is a set of processes distributed across a network of machines that
work as an ensemble to solve a common problem. In the past most were \textbf{client-server}
(resource management centralised at the server). \textbf{Peer-to-peer (P2P)} computing
represents a movement toward ``truly'' distributed applications.

\begin{figure}[H]
\centering
\begin{tikzpicture}
  % client-server
  \node[dark] (s) at (0,0) {Server};
  \foreach \i/\ang in {1/150,2/180,3/210,4/30,5/0,6/330}{
    \node[box,minimum width=1.2cm] (c\i) at ($(s)+(\ang:2.2cm)$) {\scriptsize Client};
    \draw[darr] (s)--(c\i);
  }
  \node at (0,-2.6) {\small (a) Clients invoke servers};
  % p2p
  \foreach \i/\ang in {1/90,2/162,3/234,4/306,5/18}{
    \node[box,minimum width=1.2cm] (p\i) at ($(8,0)+(\ang:1.6cm)$) {\scriptsize Peer};
  }
  \foreach \i/\j in {1/2,2/3,3/4,4/5,5/1,1/3,2/4,3/5,4/1,5/2}{
    \draw[darr] (p\i)--(p\j);
  }
  \node at (8,-2.6) {\small (b) Peer processes};
\end{tikzpicture}
\caption{Client-server versus peer-to-peer distributed applications.}
\end{figure}

\section{Grid Computing}
\begin{defbox}{Grid Computing}
\begin{itemize}[leftmargin=*]
  \item \textbf{Pcwebopedia:} harnesses \emph{unused processing cycles} of all computers in a
        network for problems too intensive for any stand-alone machine.
  \item \textbf{IBM:} enables the \emph{virtualization} of distributed computing and data
        resources (processing, network bandwidth, storage) to create a \emph{single system
        image}; a grid user sees one large virtual computer.
  \item \textbf{Sun Microsystems:} infrastructure providing \emph{dependable, consistent,
        pervasive and inexpensive} access to computational capabilities.
\end{itemize}
\end{defbox}

\textbf{Electrical power grid analogy.} Users plug into a wall socket without caring where the
electricity is generated; the power grid links many kinds of power plants. Likewise, grid users
get processors, storage, data and applications as needed with little or no knowledge of where
resources are or what hardware/OS lies beneath; ``the Grid'' links PCs, workstations, servers and
storage elements and provides the mechanism to access them.

\textbf{Features of grids:} (1) share more than information -- data, computing power and
applications in a dynamic, multi-institutional \emph{virtual organisation}; (2) efficient use of
resources at many institutes; (3) join local communities; (4) interaction with the underlying
layers must be \emph{transparent and seamless}.

\textbf{Need:} today's science is based on computation, data analysis, visualization and
collaboration; simulation and modelling are more cost effective than experiments; problems are
increasingly complex and need accurate solutions quickly; data visualization is important; and
under-utilised resources can be exploited.

\subsection{Types of Grids}
\begin{center}
\begin{tabularx}{\textwidth}{>{\bfseries}lX}
\toprule
Computational grid & Secure access to a huge pool of shared processing power for
high-throughput and computation-intensive applications.\\
Data grid & Infrastructure for storage, discovery, handling, publication and manipulation of
large data volumes stored in heterogeneous databases and file systems.\\
Collaboration grid & Supports collaboration, e.g.\ people from different companies in a virtual
enterprise working on components of a CAD project without disclosing proprietary technology.\\
Network grid & Fault-tolerant, high-performance communication; each node acts as a data router
providing caching to speed up communication.\\
Utility grid & Ultimate form: not only data and cycles but software and special equipment are
shared; applications run on one machine and users send data and receive results.\\
\bottomrule
\end{tabularx}
\end{center}

\section{Cluster Computing}
\begin{defbox}{Cluster}
A cluster is a type of parallel or distributed computer system consisting of a collection of
\textbf{inter-connected stand-alone computers} working together as a \textbf{single integrated
computing resource}. Key components: multiple standalone computers (PCs, workstations or SMPs),
operating systems, high-performance interconnects, middleware, parallel programming environments
and applications.
\end{defbox}
Clusters improve speed and/or reliability over a single computer and are more cost effective than
a single computer of comparable speed or reliability. In a typical cluster the network is faster
and closer than a LAN, uses low-latency protocols, and is \emph{more loosely coupled than SMP}.

\textbf{Types:} High Availability (failover) clusters; Load-Balancing clusters; Parallel/Distributed
processing clusters. \textbf{Components:} cluster nodes, cluster network, network characterisation.

\textbf{Benefits:} system availability (redundant hardware, OS and applications); hardware fault
tolerance (e.g.\ disk RAID); OS and application reliability; scalability (add servers or clusters);
high performance for cluster-enabled programs.

\begin{notebox}{Grid versus Cluster (Assignment 1, Q10)}
A \textbf{grid} is a form of distributed computing in which a large number of geographically
dispersed, often heterogeneous resources across multiple administrative domains are
\emph{loosely coupled} to work on a large-scale task. A \textbf{cluster} is tightly coupled and
co-located; a node is a single machine; a hypervisor is virtualization software.
\end{notebox}

\section{Utility Computing}
\begin{defbox}{Utility Computing}
Utility computing is a \emph{service provisioning model} in which a provider makes computing
resources and infrastructure management available to the customer as needed and charges for
\textbf{specific usage rather than a flat rate}. There is low or no initial cost; resources are
essentially rented. The analogy is electricity: fluctuating customer demand, charged by usage
(\emph{pay-per-use}). Utility computing is purely a concept which cloud computing practically
implements, usually with virtualization so that available storage/computing power is far larger
than that of a single time-sharing computer.
\end{defbox}

\textbf{Key elements:} (a) pay-for-use pricing model, (b) data-centre virtualization and
provisioning, (c) solves the resource-utilization problem, (d) outsourcing, (e) web services
delivery, (f) automation.

\textbf{Payment models:} flat rate, tiered, subscription, metered, pay-as-you-go, standing charges;
different customers get different models depending on scale, commitment and payment frequency; the
price expresses the provider's cost plus a profit margin.

\textbf{Risks:} data backup, data security, partner competency, defining the SLA, getting value
from charge-back.

\section{Cloud Computing: NIST Definition}
\begin{defbox}{Cloud Computing (NIST)}
Cloud computing is a model for enabling \textbf{ubiquitous, convenient, on-demand network access}
to a \textbf{shared pool of configurable computing resources} (e.g.\ networks, servers, storage,
applications and services) that can be \textbf{rapidly provisioned and released with minimal
management effort or service provider interaction}.
\end{defbox}

\subsection{Five Essential Characteristics}
\begin{enumerate}[leftmargin=*]
  \item \textbf{On-demand self-service:} a consumer can unilaterally provision capabilities such
        as server time and network storage automatically, without human interaction with each
        provider.
  \item \textbf{Broad network access:} capabilities are available over the network and accessed
        by standard mechanisms that promote use by heterogeneous thin or thick clients (mobile
        phones, tablets, laptops, workstations).
  \item \textbf{Resource pooling:} the provider's resources are pooled to serve many consumers
        using a \emph{multi-tenant model}, with physical and virtual resources dynamically
        assigned and reassigned according to demand.
  \item \textbf{Rapid elasticity:} capabilities can be elastically provisioned and released,
        sometimes automatically, to scale outward and inward with demand; to the consumer they
        often appear unlimited.
  \item \textbf{Measured service:} resource use is automatically controlled and optimised using a
        metering capability (storage, processing, bandwidth, active user accounts) and can be
        monitored, controlled and reported, giving transparency to provider and consumer.
\end{enumerate}

\subsection{Common Characteristics}
Massive scale, resilient computing, homogeneity, geographic distribution, virtualization,
service orientation, low-cost software, advanced security.

\begin{solved}{Identifying the characteristic (Assignment 1, Q1)}
\Given Statement: ``Computing resources are dynamically assigned to multiple customers using a
shared infrastructure.''
\Formula Match the statement to a NIST essential characteristic.
\Steps
\begin{enumerate}[leftmargin=*]
  \item Many customers share the same physical hardware $\Rightarrow$ multi-tenancy.
  \item Assignment and reassignment follows demand $\Rightarrow$ dynamic.
  \item These two properties together define \emph{resource pooling}.
\end{enumerate}
\Answer Resource pooling (option B in the assignment).
\end{solved}

\section{Cloud Service Models}
\begin{center}
\begin{tabularx}{\textwidth}{>{\bfseries}lXX}
\toprule
Model & Capability provided to the consumer & Consumer controls\\
\midrule
SaaS & Use provider's applications on cloud infrastructure via thin client (web browser) or
program interface. Example: Google Spreadsheet. & Nothing underneath; at most limited user-specific
configuration.\\
PaaS & Deploy consumer-created or acquired applications using languages, libraries, services and
tools supported by the provider. Examples: Microsoft Azure, Google App Engine. & Deployed
applications and possibly hosting-environment configuration; not network, servers, OS, storage.\\
IaaS & Provision processing, storage, networks and other fundamental resources; consumer deploys
arbitrary software. Examples: Amazon Web Services, Flexiscale, EC2. & OS, storage, deployed
applications and possibly some networking components.\\
\bottomrule
\end{tabularx}
\end{center}

\begin{figure}[H]
\centering
\begin{tikzpicture}[x=1cm,y=1cm]
  \foreach \kk/\xx/\nm in {9/0/Classical,5/3.2/IaaS,2/6.4/PaaS,0/9.6/SaaS}{
    \node[font=\bfseries\small,text=navy] at (\xx+1.5,0.35) {\nm};
    \foreach \rw/\lay in {1/Applications,2/Data,3/Runtime,4/Middleware,5/OS,6/Virtualization,7/Servers,8/Storage,9/Network}{
      \ifnum\rw>\kk
        \draw[fill=slate!30,draw=slate] (\xx,{-\rw*0.6}) rectangle ++(3,0.6);
      \else
        \draw[fill=navy!25,draw=navy] (\xx,{-\rw*0.6}) rectangle ++(3,0.6);
      \fi
      \node[font=\scriptsize] at (\xx+1.5,{-\rw*0.6+0.3}) {\lay};
    }
  }
  \fill[navy!25,draw=navy] (0,-6.4) rectangle ++(0.4,0.3);
  \node[anchor=west,font=\scriptsize] at (0.5,-6.25) {Managed by user};
  \fill[slate!30,draw=slate] (4,-6.4) rectangle ++(0.4,0.3);
  \node[anchor=west,font=\scriptsize] at (4.5,-6.25) {Managed by service provider};
\end{tikzpicture}
\caption{Responsibility split in the classical model, IaaS, PaaS and SaaS.}
\end{figure}

\begin{solved}{Classify the service model (Assignment 1, Q2)}
\Given A startup rents Amazon EC2 virtual servers; the team installs and patches its own OS,
configures the web server, manages runtime and applications. The provider supplies only
virtualized compute, storage and networking.
\Formula Service model = determined by the lowest layer the provider manages and the first layer the
consumer manages.
\Steps
\begin{enumerate}[leftmargin=*]
  \item Provider manages: virtualization, servers, storage, network.
  \item Consumer manages: OS, middleware/runtime, applications, data.
  \item This matches the right-most boundary of the IaaS column in the figure above.
\end{enumerate}
\Answer \textbf{IaaS}. SaaS (no consumer control), PaaS (OS is provider-managed) and FaaS are rejected.
\end{solved}

\subsection{Simplified Descriptions and Analogy}
SaaS applications are designed for end users and delivered over the web; PaaS is the set of tools
and services that makes coding and deploying applications quick and efficient; IaaS is the hardware
and software that powers it all -- servers, storage, networks, operating systems. In the
\textbf{transportation analogy}, roads are the infrastructure (IaaS), cars and trucks are the platform
(PaaS) and the people and goods are software and information (SaaS).

\section{Deployment Models}
\begin{itemize}[leftmargin=*]
  \item \textbf{Private cloud:} operated solely for one organization (e.g.\ Windows Server Hyper-V).
  \item \textbf{Community cloud:} shared by several organizations supporting a specific community goal.
  \item \textbf{Public cloud:} made available to the general public (e.g.\ Google Docs, Spreadsheet).
  \item \textbf{Hybrid cloud:} composition of two or more clouds (private, community or public),
        e.g.\ \emph{cloud bursting} for load balancing between clouds.
\end{itemize}

\section{Cloud and Virtualization}
\textbf{Virtual workspaces} are abstractions of an execution environment, dynamically available to
authorised clients through well-defined protocols, with a resource quota (CPU, memory share) and a
software configuration (OS). They are implemented on \textbf{virtual machines (VMs)}: an abstraction
of a physical host machine. A \textbf{hypervisor} (VMM) intercepts and emulates VM instructions and
manages VMs (VMware, Xen, KVM). Para-virtualization (e.g.\ Xen) approaches raw physical performance.

\begin{figure}[H]
\centering
\begin{tikzpicture}
  \node[dark,minimum width=9cm] (hw) at (0,0) {Hardware};
  \node[mid,minimum width=9cm,above=2mm of hw] (hv) {Hypervisor / Virtual Machine Monitor (Xen, VMware, KVM)};
  \node[box,minimum width=2.8cm,above=2mm of hv,xshift=-3cm] (g1) {Guest OS\\(Linux)};
  \node[box,minimum width=2.8cm,above=2mm of hv] (g2) {Guest OS\\(NetBSD)};
  \node[box,minimum width=2.8cm,above=2mm of hv,xshift=3cm] (g3) {Guest OS\\(Windows)};
  \foreach \g in {g1,g2,g3}{
    \node[gold,minimum width=1.2cm,above=2mm of \g.north west,anchor=south west] {\scriptsize App};
    \node[gold,minimum width=1.2cm,above=2mm of \g.north east,anchor=south east] {\scriptsize App};
  }
\end{tikzpicture}
\caption{Virtualized stack: several guest operating systems share one physical machine.}
\end{figure}

\textbf{Advantages of VMs:} run OSs where physical hardware is unavailable; easy to create and back up
machines; software testing on clean installs; emulate more machines than physically available;
timeshare lightly loaded systems on one host; debug by suspend/resume; easy migration; run legacy
systems.

\begin{solved}{Virtualization statements (Assignment 1, Q4 and Q7)}
\Given Four statements about VMs and hypervisors.
\Formula A VM provides a virtual machine with its own guest OS; Type-1 hypervisors run on bare metal,
Type-2 run as applications on a host OS.
\Steps
\begin{enumerate}[leftmargin=*]
  \item ``Each VM always shares the host kernel'' describes \emph{containers}, not VMs -- false.
  \item Multiple VMs can run different OSs simultaneously on one machine -- \textbf{true}.
  \item Virtualization pre-dates cloud, so ``only possible in cloud'' is false.
  \item Overcommitment (CPU/memory) is a commonly used feature, not something virtualization prevents.
  \item Type-2 (hosted) hypervisors (VirtualBox, VMware Workstation) run as an application on a host OS.
\end{enumerate}
\Answer VMs of different OSs coexist (Q4: B); Type-2 hypervisor runs on top of a host OS (Q7: B).
\end{solved}

\section{Cloud-sourcing, Cloud Storage, Advantages and Disadvantages}
\textbf{Why cloud-sourcing matters:} high-scale low-cost providers, any time/place access via browser,
rapid scalability with incremental cost and load sharing, and no need to focus on local IT.
\textbf{Concerns:} performance, reliability and SLAs; control of data and service parameters;
application features and choices; interaction between providers; no standard API (mix of SOAP and REST);
privacy, security, compliance, trust.
\textbf{Cloud storage:} large web companies rent out spare storage (e.g.\ Amazon EC2 and S3).

\subsection{Advantages}
Lower computer costs; improved performance; reduced software costs; instant software updates; improved
document-format compatibility; unlimited storage capacity (hundreds of petabytes versus a 1~TB drive);
increased data reliability (e.g.\ Dropbox, SkyDrive); universal information access; latest-version
availability; easier group collaboration; device independence.

\subsection{Disadvantages}
Requires a constant Internet connection; does not work well with low-speed connections; features may be
limited (e.g.\ PowerPoint versus web presentation tools); can be slow; stored data might not be secure;
stored data can be lost; HPC applications using MPI/OpenMP are difficult (co-location of VMs needed to
minimise latency); general concerns -- different protocols and APIs per cloud, proprietary database and
workflow systems requiring application adaptation.

\section{Business Drivers and Cloud Economics (Introductory)}
Public clouds can significantly reduce IT cost. To the end user, the cloud gives the illusion of
potentially infinite capacity with rapid scale and pay-per-use. In contrast, private data centres must be
\emph{provisioned for peak capacity}, giving a low average utilization of \textbf{5--20\%}.

\begin{center}
\begin{tabular}{lrr}
\toprule
 & In-house server & Cloud server\\
\midrule
Purchase cost & \$9600 & 0\\
Cost/hr (over 3 years) & \$0.36 & \$0.68\\
Cost ratio Cloud/In-house & \multicolumn{2}{c}{1.88}\\
Efficiency & 40\% & 80\%\\
Cost per effective hour & \$0.90 & \$0.85\\
Power and cooling & \$0.36 & 0\\
Management cost & \$0.10 & \$0.01\\
\textbf{Total cost/effective hour} & \textbf{\$1.36} & \textbf{\$0.86}\\
Cost ratio In-house/Cloud & \multicolumn{2}{c}{1.58}\\
\bottomrule
\end{tabular}
\end{center}
The in-house server in the table is an x86 machine with three quad-core processors, 12~GB RAM and a 300~GB disk.

\begin{solved}{IaaS economics: why the dearer cloud server is cheaper per useful hour}
\Given In-house: \$0.36/hr, efficiency 40\%, power and cooling \$0.36, management \$0.10.
Cloud: \$0.68/hr, efficiency 80\%, power and cooling \$0, management \$0.01.
\Formula Cost per effective hour $= \dfrac{\text{cost per hour}}{\text{efficiency}}$;\quad
total $=$ cost per effective hour $+$ power and cooling $+$ management.
\Steps
\begin{enumerate}[leftmargin=*]
  \item In-house: $0.36/0.40 = 0.90$; total $= 0.90 + 0.36 + 0.10 = 1.36$.
  \item Cloud: $0.68/0.80 = 0.85$; total $= 0.85 + 0 + 0.01 = 0.86$.
  \item Raw price ratio: $0.68/0.36 = 1.88$ (cloud dearer per hour).
  \item Total ratio: $1.36/0.86 = 1.58$.
\end{enumerate}
\Answer Total cost per effective hour: in-house \$1.36, cloud \$0.86, ratio $\approx 1.58$ in favour of the cloud.
\end{solved}

\textbf{Benefits for end users of public cloud:} high utilization and scalability; no hardware
procurement; no power cost; no IT administration; user-friendly SLAs (about 99\% availability with
compensation on a miss); rent to prototype before investing; portability standards (Cloud Computing
Interoperability Forum, Open Cloud Consortium; Amazon EC2 and Eucalyptus share the same API);
startups avoid capital investment. \textbf{Benefits for SMBs (fewer than 250 employees)} follow the same
lines. \textbf{Benefits of private cloud:} one 12-core 12~GB server costs far less than twelve 1-core 1~GB
servers; data confidentiality preserved; VMs cheaper and faster to provision than physical machines.

\textbf{Economics of PaaS versus IaaS:} for a 24x7 web application with unpredictable volume, IaaS requires
a minimum set of servers always provisioned, while deploying on PaaS costs nothing at rest, costs are
incurred by usage, and the platform scales automatically. \textbf{PaaS benefits:} no scaling or load
balancing by the user; web-based IDE; easier migration from development to production; no separate IDE
licences. \textbf{SaaS benefits:} subscribe instead of buying and licensing software; use free tools such as
Google Docs; enterprises use web-based CRM instead of installing servers.

\textbf{Purposes of cloud in organizations:} IT platform for multi-organization processes; backup; CRM,
ERP, supply-chain; personal productivity and collaboration; development and test; storing and archiving
large files; analysing customer/operations data; e-business and e-government sites; R\&D analysis; handling
demand spikes; batch processing (conversion, risk modelling, rendering); sharing information with
regulators; consumer entertainment. \textbf{Top cloud applications:} mail and messaging, archiving, backup,
storage, security, virtual servers, CRM, cross-enterprise collaboration, hosted PBX, video conferencing.

\section{Cloud Computing Architecture}
\subsection{Building Blocks}
\begin{itemize}[leftmargin=*]
  \item \textbf{Technical architecture:} structuring by the XaaS stack; adopting cloud paradigms;
        structuring services and components; relationships and external endpoints; middleware and
        communication; management and security.
  \item \textbf{Deployment/operation architecture:} geo-location check (legal issues, export control);
        operation and monitoring.
\end{itemize}

\subsection{XaaS}
\begin{defbox}{XaaS}
Combining Service-Oriented Infrastructure (SOI) with cloud computing realises \textbf{XaaS}: ``anything'' or
``everything'' as a Service -- a growing number of services delivered over the Internet rather than
locally. \textbf{Broad network access + resource pooling (cloud) + business-driven infrastructure on demand +
service orientation (SOI) = XaaS}, which fulfils the CSP demands: increase productivity, end-user
satisfaction, innovation and agility.
\end{defbox}
Other XaaS examples: Business Process as a Service (BPaaS), Storage as a Service, Security as a Service
(SECaaS), Database as a Service (DaaS), Monitoring/Management as a Service (MaaS), Communications, Content and
Computing as a Service (CaaS), Identity as a Service (IDaaS), Backup as a Service (BaaS), Desktop as a Service.

\begin{notebox}{Assignment 1, Q6}
The combination of SOA and cloud computing leading to ``almost any IT capability delivered as a service'' is
called \textbf{XaaS}. FaaS (Functions as a Service) is just one specific model.
\end{notebox}

\subsection{Classical (Legacy) Model versus Evergreen IT}
In the classical model all layers (hardware, OS, development tools, applications) are managed by users who
bear the hardware, maintenance and technology cost; systems are custom build-to-order, deployed as
tightly coupled vertical stacks of layers (no single part easily replaced), with manual provisioning --
\emph{Legacy IT}. Cloud simplifies the stack towards \emph{Evergreen IT}: dedicated infrastructure,
application and end-user devices form a simplified IT stack.

\subsection{Client-Server Model versus Cloud Model}
\begin{center}
\begin{tabularx}{\textwidth}{>{\bfseries}XXX}
\toprule
 & Client-server model & Cloud computing model\\
\midrule
Service model & Simple: server services client requests & Variety of complex models (IaaS, PaaS, SaaS)\\
Load balancing & May or may not be load balanced & Load balanced\\
Scalability & Scalable to some extent in a cluster & Theoretically infinitely scalable\\
Virtualization & No concept of virtualization & Virtualization is the core concept\\
\bottomrule
\end{tabularx}
\end{center}
A three-tier client-server architecture has presentation, application and data tiers; the client is often a
thin client, and the server may be a load-balanced cluster or a stand-alone machine.

\subsection{SaaS, PaaS, IaaS in Detail}
\begin{center}
\begin{tabularx}{\textwidth}{>{\bfseries}lXXX}
\toprule
 & SaaS & PaaS & IaaS\\
\midrule
Definition & Software deployed over the Internet, licensed on demand, by subscription, pay-as-you-go, or free
(ad/user-list revenue). & Platform for creating web applications quickly without buying and maintaining
underlying software and infrastructure. & Servers, storage, network and OS delivered as an on-demand,
outsourced service.\\
Characteristics & Web access to commercial software; managed from a central location; one-to-many delivery;
no upgrades or patches by users; APIs for integration. & Integrated development, test, deploy, host and
maintain services; web-based UI tools; multi-tenant; built-in scalability, load balancing and failover;
web-service/database integration; team collaboration; billing and subscription tools. & Resources
distributed as a service; dynamic scaling; variable utility pricing; multiple users per hardware.\\
Good fit & Interplay with the outside world (e.g.\ email campaigns); web/mobile access; short-term use;
demand spikes (tax/billing); e.g.\ Salesforce CRM. & Multiple developers or external parties; automated
test and deployment; agile development; e.g.\ Azure, Google App Engine. & Volatile demand; new
organizations without capital; rapid growth; limiting CAPEX in favour of OPEX; trial or temporary needs.\\
Poor fit & Extremely fast real-time processing; regulations forbidding external hosting; existing on-premise
solution suffices. & Highly portable applications; proprietary languages causing lock-in; need to
customise underlying hardware/software. & Regulatory limits on offshoring/outsourcing data; highest
performance needs met by on-premise hardware.\\
\bottomrule
\end{tabularx}
\end{center}

\begin{solved}{Advantages of PaaS (Assignment 1, Q8)}
\Given Four claimed advantages: (i) developers focus on application development, (ii) automatic software
updates of the underlying platform, (iii) complete control over physical server hardware, (iv) faster
application deployment.
\Formula PaaS abstracts infrastructure management and offers managed runtime environments.
\Steps
\begin{enumerate}[leftmargin=*]
  \item (i) true -- infrastructure is abstracted. (ii) true -- the provider updates the platform.
  \item (iii) false -- physical hardware is the provider's responsibility.
  \item (iv) true -- ready-made runtime and IDE speed deployment.
\end{enumerate}
\Answer (i), (ii) and (iv) -- option A.
\end{solved}

\section{Role of Networking and Network Function Virtualization}
Network resources can be provisioned dynamically in the cloud. Core concepts: Virtual LANs, Virtual
Private Networks and protocol layers; tools such as OpenSSH and OpenVPN set up topologies and
configurations.

\begin{defbox}{Network Functions Virtualisation (ETSI NFV)}
NFV aims to transform how operators architect networks by evolving standard IT virtualisation technology to
consolidate many network-equipment types onto industry-standard high-volume servers, switches and storage in
data centres, network nodes and end-user premises. Network functions are implemented in software that can be
moved to, or instantiated in, various locations as required, without installing new equipment.
\end{defbox}

\section{Week 1 Assignment Review}
\begin{center}
\begin{tabularx}{\textwidth}{cXc}
\toprule
Q & Concept tested & Answer\\
\midrule
1 & Resource pooling (multi-tenant shared resources) & B\\
2 & EC2 rented VM with user-managed OS $\Rightarrow$ IaaS & C\\
3 & Private cloud: exclusive use, on-premises or third-party hosted & D\\
4 & Different OS VMs simultaneously on one machine & B\\
5 & Distributed systems: scalability and fault tolerance, \emph{not} guaranteed lower latency & A\\
6 & SOA + cloud $\Rightarrow$ XaaS & C\\
7 & Type-2 hypervisor runs as an application on a host OS & B\\
8 & PaaS advantages (i), (ii), (iv) & A\\
9 & Amazon EC2 is IaaS & B\\
10 & Grid: geographically dispersed, heterogeneous, loosely coupled & B\\
\bottomrule
\end{tabularx}
\end{center}

\begin{solved}{Distributed system statements (Assignment 1, Q5)}
\Given (i) adding nodes improves scalability; (ii) replication increases fault tolerance;
(iii) distribution guarantees lower latency for every operation than a centralized system.
\Formula Distribution adds network communication and coordination overhead.
\Steps
\begin{enumerate}[leftmargin=*]
  \item (i) True: horizontal scaling.
  \item (ii) True: replication of data and services.
  \item (iii) False: remote calls and coordination can make some operations slower.
\end{enumerate}
\Answer (i) and (ii) only -- option A.
\end{solved}
